This thesis has investigated the use of a Basilisk simulation framework for estimating the fuel-limited lifetime of a simplified \gls{gnssr} formation flying mission. The work was divided into two main parts. First, the long-horizon translational behavior of the selected Basilisk propagation configuration was benchmarked against a continuously updating Skyfield-\gls{sgp4} reference. Second, the simulator was expanded into a \gls{gnssr} formation flight mission simulator, including spacecraft subsystems, operational mode switching, attitude control, formation control, propulsion, power consumption, and fuel depletion.

The first research question asked how the final high-fidelity Basilisk propagation configuration compares against a continuously updating Skyfield-\gls{sgp4} reference over a long-horizon propagation interval. The comparative propagation analysis showed that the primary thesis Basilisk configuration, using the NRLMSISE-00 atmosphere model, diverged less from the updating-\gls{tle} Skyfield-\gls{sgp4} baseline than both exponential-atmosphere Basilisk configurations. The specialization-project exponential atmosphere configuration produced the largest same-spacecraft state difference, while the re-tuned exponential atmosphere case reduced the divergence but did not outperform the NRLMSISE-00 configuration.

These results indicate that the thesis Basilisk configuration is the most suitable of the evaluated configurations for the subsequent formation flight lifetime analysis. However, the comparison does not constitute a strict validation of absolute orbit prediction accuracy, since the Skyfield-\gls{sgp4} baseline should be treated as a self-correcting reference rather than a true orbit-determination ground truth. The propagation benchmark therefore provides confidence and uncertainty context for the lifetime analysis, rather than exact long-horizon prediction error bounds.


\begin{comment}
* The second research question asked what fuel-limited mission lifetime is predicted for the selected \gls{gnssr} leader-follower formation when formation acquisition, steady-state formation maintenance, spacecraft subsystem constraints, and operational mode switching are included in the simulation.
* 2 test campaigns was performed in order to fit 2 models for acquisition and steady-state fuel consumption, respectively. 
* Using this model to predict the lifetime of the gnss-r formation flight mission, 
\end{comment}


The second research question asked what fuel-limited mission lifetime is predicted for the selected \gls{gnssr} leader-follower formation when formation acquisition, steady-state formation maintenance, spacecraft subsystem constraints, and operational mode switching are included in the simulation. To answer this, two separate simulation campaigns were performed. The deployment velocity sensitivity campaign was used to estimate the acquisition fuel consumption as a function of differential deployment speed, while the steady-state fuel consumption campaign was used to estimate the post-acquisition formation-maintenance fuel consumption over time.

In the final lifetime estimate, the acquisition and steady-state fuel consumption models were combined into a remaining-fuel model. This model was then used to predict the fuel-limited lifetime of the \gls{gnssr} formation flight mission for the differential deployment speed cases considered in the deployment sensitivity campaign. The resulting linear lifetime estimates ranged from $413.6~\mathrm{days}$ for the most optimistic zero-differential-deployment case to $366.5~\mathrm{days}$ for the $\Delta v_{\mathrm{dep}}=4~\mathrm{m/s}$ case.

The acquisition fuel model contains uncertainty related to the definition of acquisition time. The extracted acquisition fuel depends on the selected orbital element difference tolerances and dwell-time criterion, and some post-acquisition burns still contributed noticeably to the cumulative fuel consumption. However, the steady-state fuel consumption model is the dominant source of lifetime estimation uncertainty. This is because the predicted lifetime is extrapolated over several hundred days, while the completed steady-state campaign lasted only $24$ days. The linear model used for the final lifetime estimate also does not capture the observed reduction in fuel consumption rate over the simulated interval. Since a decreasing stored propellant mass reduces the propellant required for equivalent orbital corrections, this makes the linear lifetime estimate likely pessimistic with respect to this effect.

The comparative propagation analysis gives a second indication that the lifetime estimate may be pessimistic. The same-formation comparison showed larger leader-follower drift in the Basilisk configurations than in the updating-\gls{tle} Skyfield-\gls{sgp4} baseline. Since formation-maintenance fuel is driven by the correction of relative drift, an overestimated drift rate would also lead to overestimated steady-state fuel consumption. However, this interpretation is limited by the fact that the Skyfield-\gls{sgp4} baseline is not a true physical ground truth, and by the fact that the same-formation comparison was performed in a deactivated open-loop propagation scenario rather than in the closed-loop formation-maintenance simulator.

The main limitations of the thesis are tied to the considered formation scenario, the available simulation duration, and the implemented formation controller. The lifetime analysis considered only a two-spacecraft leader-follower concentric formation. The resulting lifetime estimate is therefore specific to the selected formation geometry and cannot be generalized directly to other cooperative \gls{gnssr} formations, larger satellite groups, or additional followers in the concentric formation. More simulation campaigns would be required to determine how the acquisition and maintenance fuel consumption change for other formation types, other radii, and more spacecraft.

The simulation duration is another major limitation. The deployment sensitivity campaign was limited to $48~\mathrm{h}$ per case, while the steady-state campaign was intended for a longer horizon but produced a final completed run of only $24$ days due to unresolved memory overflow issues. The final lifetime estimate therefore depends on extrapolating a fitted maintenance fuel model far beyond the simulated interval. Longer completed steady-state simulations would be required to determine whether the observed reduction in fuel consumption rate continues, stabilizes, or changes over longer mission durations.

Finally, the achieved formation accuracy limits how directly the simulated mission can be interpreted as a complete cooperative \gls{gnssr} imaging mission. The implemented formation controller is based on classical orbital-element differences and was able to maintain a bounded relative trajectory, but it did not achieve the approximately $10~\mathrm{m}$ relative-position accuracy associated with the close-formation cooperative processing benefits discussed for HydroSwarm-like concentric formations. Future work should therefore investigate alternative control formulations, such as relative-state tracking, relative orbital elements better suited for near-circular orbits, or direct \gls{rtn} tracking. Further work should also include realistic navigation uncertainty, longer-duration simulation capability, larger formation scenarios, and explicit \gls{gnssr} observation opportunity modelling. 


% The second research question asked what fuel-limited mission lifetime is predicted for the selected \gls{gnssr} leader-follower formation when formation acquisition, steady-state formation maintenance, spacecraft subsystem constraints, and operational mode switching are included in the simulation. The deployment velocity sensitivity campaign showed that all tested deployment cases acquired the desired formation within the 48-hour simulation interval. The acquisition fuel consumption generally increased with differential deployment speed, from $5.60~\mathrm{g}$ for the zero-deployment-speed case to $77.02~\mathrm{g}$ for the $\Delta v_{\mathrm{dep}} = 4.0~\mathrm{m/s}$ case. Higher deployment speeds also generally required more burns and longer average thruster on-times.

% The steady-state fuel consumption campaign was used to estimate the post-acquisition maintenance fuel model. The completed long-horizon simulation lasted $24$ days, and the cumulative follower fuel consumption was fitted using both linear and quadratic models. The linear model provided a usable finite lifetime estimate, while the quadratic model did not reach zero fuel when extrapolated and therefore produced an undefined fuel-depletion time. Combining the fitted acquisition and maintenance fuel models produced predicted fuel-limited lifetimes from $413.6~\mathrm{days}$ for the $\Delta v_{\mathrm{dep}} = 0~\mathrm{m/s}$ case to $366.5~\mathrm{days}$ for the $\Delta v_{\mathrm{dep}} = 4.0~\mathrm{m/s}$ case. These values should be interpreted as simulation-based fuel-limited lifetime estimates for the selected spacecraft, controller, formation geometry, and disturbance environment.

% The implemented \gls{gnssr} formation flight simulator captured several mission-relevant constraints that are often absent from idealized formation-flying studies. The spacecraft model included representative mass, power, propulsion, reaction-wheel, communication, and payload power components. The operational mode logic enforced transitions between charging, communication, capture, coasting, and burn-related modes, and no burn requests were denied in the reported campaigns. This indicates that the selected mission scenario remained feasible under the implemented power and operational constraints. At the same time, the simulator did not model actual \gls{gnssr} measurement opportunities, delay-Doppler processing, cooperative observation scheduling, navigation uncertainty, hardware degradation, or a full multi-spacecraft HydroSwarm-like formation.

% The most important limitation of the final lifetime estimate is that it extrapolates a steady-state fuel model over several hundred days from a completed simulation interval of only $24$ days. The result is therefore sensitive to the assumed long-term maintenance fuel rate. The acquisition fuel model also has limitations, particularly at low deployment speeds, where the unconstrained linear fit predicts a physically unrealistic negative fuel consumption at zero deployment speed. In addition, the achieved formation accuracy does not demonstrate the approximately metre- to ten-metre-level relative-position accuracy associated with close cooperative \gls{gnssr} imaging concepts. The estimated lifetime should therefore be interpreted as a fuel-budget feasibility indicator rather than a mission-certified lifetime prediction.

% Future work should first improve the propagation benchmark by comparing the Basilisk trajectory against a precise orbit-determination reference or onboard navigation data, if available. The GNSS-R Formation Flight Simulator should also be improved to support longer simulations without memory overflow, allowing the steady-state fuel model to be fitted over a significantly longer interval. Further work should investigate alternative formation-control strategies capable of tighter relative-position control, include realistic navigation and actuation errors, and extend the analysis from a two-spacecraft leader--follower case to larger cooperative formations. Finally, the operational model should be expanded to include GNSS-R observation geometry, data storage, downlink scheduling, and cooperative capture constraints, so that future lifetime estimates can be connected more directly to science-producing mission operations.

% Overall, this thesis demonstrates that a Basilisk-based end-to-end simulator can be used to connect orbit propagation, spacecraft subsystems, formation control, and fuel consumption in a single GNSS-R formation flight lifetime analysis. The final results suggest that the selected simplified formation can remain fuel-capable for approximately one year under the simulated assumptions. The work also identifies the modelling improvements required before such simulations can be used as a higher-confidence design tool for realistic cooperative \gls{gnssr} missions.
